\documentclass[aspectratio=169]{beamer}

% Shared Beamer settings for Elements of AIML lectures (UPES).
% Keep this file free of \documentclass to allow reuse via \input.

% Allow lecture sources to use shared theme/assets without copying them into each lecture folder.
% NOTE: These paths are relative to `Unit-xx/Lecture-yy/latex/` (and `Lab/Experiment-xx/latex/`).
\makeatletter
\def\input@path{{../../../_shared/}{../../../_shared/UPES-Beamer-Theme/}}
\makeatother

\usetheme{UPES}
\setbeamertemplate{navigation symbols}{}

\usepackage{amsmath}
\usepackage{amssymb}
\usepackage{booktabs}
\usepackage{graphicx}
\usepackage{hyperref}
\usepackage{xcolor}
\usepackage{listings}

% Code listings (general-purpose).
\lstset{
  basicstyle=\ttfamily\small,
  keywordstyle=\color{blue},
  commentstyle=\color{gray},
  stringstyle=\color{teal},
  showstringspaces=false,
  columns=fullflexible,
  breaklines=true
}

\usepackage{tikz}
\usetikzlibrary{arrows.meta,positioning,shapes.geometric,calc}

% Per-slide source line (references shown on the slide itself).
\newcommand{\slidesources}[1]{%
  \vfill
  {\tiny\color{gray}\raggedright\sloppy\parbox{\linewidth}{Sources: #1}\par}%
}

% Unit-wise source strings for this subject (used by slides/notes).
% Path is relative to the lecture `latex/` folder (the compilation CWD).
% Source strings for Elements of AIML (used by slides + notes).

\newcommand{\AIMLRepoURL}{https://github.com/tali7c/Elements-of-AIML}

\newcommand{\AIMLLabSources}{%
Provost and Fawcett, \textit{Data Science for Business}.;%
McKinney, \textit{Python for Data Analysis}.;%
WEKA Documentation (Waikato Environment for Knowledge Analysis), accessed 2026-02-28.;%
Matplotlib and Seaborn documentation, accessed 2026-02-28.%
}

\newcommand{\AIMLUnitISources}{%
Rich and Knight, \textit{Artificial Intelligence}, McGraw Hill, 2017.;%
Russell and Norvig, \textit{Artificial Intelligence: A Modern Approach} (4e), Ch. 1--2 (Introduction, Intelligent Agents).;%
Poole and Mackworth, \textit{Artificial Intelligence: Foundations of Computational Agents} (2e), selected sections.%
}

\newcommand{\AIMLUnitIISources}{%
Russell and Norvig, \textit{Artificial Intelligence: A Modern Approach} (4e), Ch. 7--9 (Logical Agents, First-Order Logic, Inference).;%
Huth and Ryan, \textit{Logic in Computer Science} (2e), selected sections on propositional and first-order logic.;%
Genesereth and Nilsson, \textit{Logical Foundations of Artificial Intelligence}, selected chapters.%
}

\newcommand{\AIMLUnitIIISources}{%
Mueller and Massaron, \textit{Machine Learning for Dummies}, For Dummies, 2016.;%
Mitchell, \textit{Machine Learning}, selected chapters on learning basics and evaluation.;%
Scikit-learn documentation, accessed 2026-02-28.%
}

\newcommand{\AIMLUnitIVSources}{%
Mueller and Massaron, \textit{Machine Learning for Dummies}, For Dummies, 2016.;%
James, Witten, Hastie, and Tibshirani, \textit{An Introduction to Statistical Learning} (2e), selected chapters.;%
Scikit-learn documentation, accessed 2026-02-28.%
}

\newcommand{\AIMLUnitVSources}{%
NIST, \textit{AI Risk Management Framework (AI RMF 1.0)}, 2023.;%
Barocas, Hardt, and Narayanan, \textit{Fairness and Machine Learning} (online book), selected sections.;%
Knaflic, \textit{Storytelling with Data}, selected chapters on communicating results.%
}



\title[Elements of AIML]{Elements of AIML}
\subtitle{Unit 02 -- Lecture 02: Predicate Logic (First-Order Logic)}
\author{Tofik Ali}
\institute{School of Computer Science, UPES Dehradun}
\date{\today}

\graphicspath{{../images/}{../../../_shared/}{../../../_shared/UPES-Beamer-Theme/}}

\begin{document}

\UPESCoverFrame
\setcounter{framenumber}{0}

\begin{frame}
  \titlepage
  \vspace{0.3em}
  \begin{center}
    \footnotesize Topic focus: \textbf{objects, relations, and quantifiers} for richer knowledge bases\\[0.25em]
    \footnotesize Repository: \texttt{\AIMLRepoURL}
  \end{center}
  \slidesources{\AIMLUnitIISources}
\end{frame}

\begin{frame}{Quick Links}
  \centering
  \hyperlink{sec:why}{\beamerbutton{Why FOL}}\hspace{0.6em}
  \hyperlink{sec:syntax}{\beamerbutton{Syntax}}\hspace{0.6em}
  \hyperlink{sec:sem}{\beamerbutton{Semantics}}\hspace{0.6em}
  \hyperlink{sec:translate}{\beamerbutton{Translate}}\hspace{0.6em}
  \hyperlink{sec:summary}{\beamerbutton{Summary}}
\end{frame}

\begin{frame}{Agenda}
  \tableofcontents
\end{frame}

\section{Why predicate logic?}
\label{sec:why}

\begin{frame}{Learning Outcomes}
  By the end of this lecture, you should be able to:
  \begin{itemize}[<+->]
    \item Explain why propositional logic is often too limited for AI knowledge bases.
    \item Write basic first-order logic (FOL) sentences with predicates and quantifiers.
    \item Interpret the meaning of $\forall$ and $\exists$ and understand scope.
    \item Translate simple English statements into FOL (and spot common pitfalls).
  \end{itemize}
\end{frame}

\begin{frame}{Abbreviations \& Symbols}
  \begin{itemize}[<+->]
    \item \textbf{FOL}: First-Order Logic (predicate logic)
    \item \textbf{PL}: Propositional Logic
    \item $\forall$: ``for all'' \hspace{1.2em} $\exists$: ``there exists''
    \item $\lnot,\land,\lor,\rightarrow$: NOT, AND, OR, IMPLIES
    \item $A(x)$: predicate $A$ applied to object $x$
  \end{itemize}
\end{frame}

\begin{frame}{Why Propositional Logic Becomes Verbose}
  Propositional logic treats facts as atomic symbols.
  \vspace{0.4em}
  \begin{exampleblock}{Example}
    ``Ali is a student'' and ``Sara is a student'' become two unrelated atoms:
    \[
      Student(Ali),\; Student(Sara)\quad \text{(FOL)} \;\;\;\; \text{vs}\;\;\;\; p,\; q\quad \text{(PL)}
    \]
  \end{exampleblock}
  With many objects, PL needs too many symbols and cannot express general rules compactly.
\end{frame}

\section{Syntax}
\label{sec:syntax}

\begin{frame}{Core Building Blocks (Syntax)}
  \begin{itemize}[<+->]
    \item \textbf{Constants}: specific objects (e.g., $Ali$, $UPES$)
    \item \textbf{Variables}: placeholders (e.g., $x, y$)
    \item \textbf{Functions}: map objects to objects (e.g., $FatherOf(x)$)
    \item \textbf{Predicates}: relations / properties (e.g., $Student(x)$, $Likes(x, AI)$)
    \item \textbf{Quantifiers}: $\forall$ (for all), $\exists$ (there exists)
  \end{itemize}
\end{frame}

\begin{frame}{Atomic Sentences and Connectives}
  \begin{itemize}[<+->]
    \item \textbf{Atomic sentence}: predicate applied to terms, e.g. $Student(Ali)$ or $Likes(Ali, AI)$.
    \item Build complex sentences using propositional connectives:
    \[
      \lnot,\; \land,\; \lor,\; \rightarrow,\; \leftrightarrow
    \]
    \item Add quantifiers:
    \[
      \forall x\; Student(x) \rightarrow WorksHard(x)
    \]
  \end{itemize}
\end{frame}

\begin{frame}{Free vs Bound Variables (Scope)}
  \begin{itemize}[<+->]
    \item A variable is \textbf{bound} if it is under a quantifier.
    \item A variable is \textbf{free} if it is not bound by any quantifier.
  \end{itemize}
  \vspace{0.4em}
  Examples:
  \[
    \forall x\; Likes(x, AI) \quad \text{($x$ bound)}
  \]
  \[
    Likes(x, AI) \quad \text{($x$ free; meaning depends on context)}
  \]
\end{frame}

\begin{frame}{Quantifiers (Meaning + Scope)}
  \begin{itemize}[<+->]
    \item $\forall x\; P(x)$: for every object $x$ in the domain, $P(x)$ is true.
    \item $\exists x\; P(x)$: there exists at least one object $x$ in the domain such that $P(x)$ is true.
    \item \textbf{Scope matters}: where the quantifier applies.
  \end{itemize}
  \vspace{0.4em}
  \begin{exampleblock}{Common confusion}
    $\forall x\; \exists y\; Likes(x,y)$ is \textbf{not} the same as $\exists y\; \forall x\; Likes(x,y)$.
  \end{exampleblock}
\end{frame}

\begin{frame}{Negating Quantifiers (Very Useful Later)}
  \begin{itemize}[<+->]
    \item $\lnot(\forall x\; P(x)) \equiv \exists x\; \lnot P(x)$
    \item $\lnot(\exists x\; P(x)) \equiv \forall x\; \lnot P(x)$
  \end{itemize}
  \vspace{0.4em}
  These rules are key in clause form conversion (Lecture 03).
\end{frame}

\section{Semantics}
\label{sec:sem}

\begin{frame}{Semantics (What Makes a Sentence True?)}
  To interpret FOL, we specify:
  \begin{itemize}[<+->]
    \item a \textbf{domain} of objects (what $x$ can be),
    \item an \textbf{interpretation} that assigns meaning to constants, functions, and predicates.
  \end{itemize}
  \vspace{0.4em}
  \begin{block}{Model idea}
    A model is an interpretation where a sentence evaluates to \textbf{true}.
  \end{block}
\end{frame}

\begin{frame}{A Tiny Semantics Example}
  Domain: $\{Ali, Sara\}$ \quad Predicate: $Student(\cdot)$
  \vspace{0.4em}
  Suppose interpretation says:
  \[
    Student(Ali)=T,\;\; Student(Sara)=F
  \]
  Then:
  \[
    \exists x\; Student(x)\; \text{is true, but}\; \forall x\; Student(x)\; \text{is false.}
  \]
\end{frame}

\section{Translation practice}
\label{sec:translate}

\begin{frame}{Translation Patterns (English $\rightarrow$ FOL)}
  \begin{itemize}[<+->]
    \item ``All A are B'':
    \[
      \forall x\; A(x) \rightarrow B(x)
    \]
    \item ``Some A is B'':
    \[
      \exists x\; A(x) \land B(x)
    \]
    \item ``No A is B'':
    \[
      \forall x\; A(x) \rightarrow \lnot B(x)
    \]
  \end{itemize}
\end{frame}

\begin{frame}{More Patterns: ``Only'', ``Exactly One'' (Optional)}
  \begin{itemize}[<+->]
    \item ``Only A are B'' means: if $B$ then $A$
    \[
      \forall x\; B(x) \rightarrow A(x)
    \]
    \item ``Exactly one A exists'' (sketch):
    \[
      \exists x\; A(x)\; \land\; \forall y\; (A(y) \rightarrow y=x)
    \]
  \end{itemize}
\end{frame}

\begin{frame}{Quantifier Order Example}
  \begin{block}{Statement 1}
    Every student is enrolled in \emph{some} course:
    \[
      \forall x\; Student(x) \rightarrow \exists y\; (Course(y) \land Enrolled(x,y))
    \]
  \end{block}
  \vspace{0.4em}
  \begin{block}{Statement 2}
    There exists a single course that \emph{all} students are enrolled in:
    \[
      \exists y\; Course(y) \land \forall x\; (Student(x) \rightarrow Enrolled(x,y))
    \]
  \end{block}
\end{frame}

\begin{frame}{Common Pitfalls}
  \begin{itemize}[<+->]
    \item Forgetting to specify the domain (what objects exist?)
    \item Using $\forall x\; A(x) \land B(x)$ when you meant implication
    \item Wrong quantifier order (changes meaning)
    \item Ambiguous English (``someone'' / ``the'' / ``only'')
  \end{itemize}
\end{frame}

\begin{frame}{Bridge to Next Lecture: Substitution \& Unification (Preview)}
  In resolution for FOL, literals may ``match'' after substituting variables.
  \vspace{0.4em}
  Example:
  \[
    Likes(x, AI) \text{ matches } Likes(Ali, AI) \text{ with substitution } \{x/Ali\}.
  \]
  We will formalize this as \textbf{unification} in Lecture 03.
\end{frame}

\section{Summary}
\label{sec:summary}

\begin{frame}{Summary}
  \begin{itemize}
    \item Predicate logic adds \textbf{objects and relations} to propositional logic.
    \item Quantifiers $\forall$ and $\exists$ express general rules compactly.
    \item Semantics requires a domain and interpretation.
    \item Next: clause form, resolution, and unification (Lecture 03).
  \end{itemize}
  \slidesources{\AIMLUnitIISources}
\end{frame}

\begin{frame}{Exit Questions}
  \begin{enumerate}
    \item Why is propositional logic insufficient for many knowledge bases?
    \item Translate: ``All humans are mortal'' and ``Some student likes AI'' into FOL.
    \item Explain why $\forall x \exists y\, P(x,y)$ differs from $\exists y \forall x\, P(x,y)$.
    \item Give one example of an English sentence that is ambiguous without clarifying assumptions.
  \end{enumerate}
\end{frame}

\end{document}
